\chapter{Python and ROS 2: nodes and topics}
\label{ch:ros2}
The algorithms developed so far must eventually exchange data with sensors,
actuators, and other algorithms. A controller needs the robot's estimated
pose and a reference to follow. A mapping algorithm needs measurements and
the poses from which they were collected. A planner needs a map, a start,
and a goal. Implementing these functions as communicating components makes
their responsibilities and interfaces explicit.

This chapter introduces enough Python to read two small ROS 2 programs:
a periodic publisher and a subscriber that processes its messages. The
examples use text so that the communication is easy to observe. The same
structure will later carry measurements and robot commands.

\section{A short introduction to Python}
\subsection{Loops, conditions, and functions}
In Python, indentation determines which statements belong to a block.
A colon introduces the body of a loop, conditional, function, or class.
Use consistent indentation; four spaces per level is the usual convention.
Unlike C braces or MATLAB's \texttt{end}, indentation is part of the syntax.

\begin{lstlisting}[style=robotcode,language=Python]
def thisIsTwooo():
    print("THIS IS TWOOOOO")

for i in range(0, 5):
    print(i)
    if i == 2:
        print("this is two!")
thisIsTwooo()
\end{lstlisting}
The loop visits $0,1,2,3,4$: the upper limit of \texttt{range} is excluded.
The operator \texttt{==} compares values; \texttt{=} assigns a value.
The conditional message appears only when \texttt{i} equals 2. The final
function call is outside the loop and therefore executes once, after it.
Defining a function with \texttt{def} does not execute its body; calling
the function does.

\begin{minipage}{\linewidth}
The output is
\begin{lstlisting}[style=robotcode]
0
1
2
this is two!
3
4
THIS IS TWOOOOO
\end{lstlisting}
\end{minipage}

Moving \texttt{thisIsTwooo()} one indentation level to the right places
it inside the loop, but outside the conditional: it then executes five
times. Moving it another level to the right places it inside the conditional,
so it executes once, immediately after \texttt{this is two!}. These programs
differ in behavior even though the words on the line are unchanged.

\subsection{Objects, methods, and instance data}
A class describes objects that combine data with operations on those data.
An \emph{instance} is a particular object created from the class. A function
defined in a class is a method; the parameter \texttt{self} refers to the
instance on which an ordinary instance method operates.
\begin{lstlisting}[style=robotcode,language=Python]
class Counter:
    def __init__(self):
        self.i = 0

    def increment(self):
        self.i += 1

counter = Counter()
counter.increment()
print(counter.i)       # 1
\end{lstlisting}
The initialization method \texttt{\_\_init\_\_} runs when the instance is
initialized. Here it gives that instance an attribute named \texttt{i}.
The method call \texttt{counter.increment()} operates on that same object.
An ordinary local variable exists within a function's execution; an
attribute such as \texttt{self.i} stores state that other methods of the
instance can access later. This is useful for a counter that persists
between successive callbacks.

Inheritance lets a new class extend an existing class. In
\texttt{class MinimalPublisher(Node)}, the new class inherits ROS-related
methods from \texttt{Node}, then defines its own initialization and callback.
Calling \texttt{super().\_\_init\_\_(...)} performs the base-class
initialization needed before using those ROS facilities.

\section{A robot as a network of software components}
\label{sec:l14-architecture}
A robot driver translates between the hardware interface and software
messages. For example, it may read wheel information over a serial link,
publish odometry, and translate requested linear and angular velocities
into hardware commands. A lidar driver similarly translates scan data
into a representation that other software can consume. USB serial links
are one possible hardware interface, not a requirement of ROS 2.

Above these drivers, separate components can implement mapping, planning,
and control, as in Figure~\ref{fig:l14-architecture}:
\begin{itemize}
\item The mapping component combines scans with an available pose estimate.
\item The planner uses the map, start, and goal to produce a path.
\item The controller combines a reference with pose feedback to compute
velocity commands $(v,\omega)$.
\end{itemize}
A geometric path still needs an execution policy: a controller may select
successive waypoints, or a separate step may assign timing to obtain a
trajectory. The diagram groups this reference-generation step with planning.

\begin{figure}[htbp]
\centering\begin{tikzpicture}[x=1cm,y=1cm,>=Stealth,
 box/.style={draw,rounded corners,align=center,minimum width=2.1cm,minimum height=1.05cm},
 flow/.style={->,thick,sensorframe},
 lab/.style={font=\footnotesize,fill=white,inner sep=2pt}]
\draw[black!40,rounded corners] (-1.45,-1.25) rectangle (9.1,5.6);
\node[anchor=west,font=\footnotesize] at (-1.3,5.35) {Computer: ROS 2 software};
\node[box,draw=robotframe] (plan) at (0,3.8) {Path planner\\and reference};
\node[box,draw=robotframe] (control) at (4,3.8) {Controller};
\node[box,draw=robotframe] (map) at (0,0) {Mapping};
\node[box] (driver) at (7.6,3.8) {Robot driver};
\node[box] (lidardriver) at (7.6,0) {Lidar driver};
\node[box,minimum width=1.6cm] (robot) at (11,3.8) {Robot};
\node[box,minimum width=1.6cm] (lidar) at (11,0) {Lidar};
\draw[flow] (map) -- node[lab,right] {map} (plan);
\draw[flow] (plan) -- node[lab,above] {reference} (control);
\draw[flow] ([yshift=-.22cm]control.east) -- node[lab,below] {$v,\omega$} ([yshift=-.22cm]driver.west);
\draw[flow] ([yshift=.22cm]driver.west) -- node[lab,above] {odometry} ([yshift=.22cm]control.east);
\draw[flow] (driver.south) -- (7.6,1.7) -- (2,1.7) -- (2,.3) -- ([yshift=.3cm]map.east);
\node[lab] at (3.8,1.7) {odometry / pose estimate};
\draw[flow] (lidardriver) -- node[lab,above] {scan} (map);
\draw[->,thick] ([yshift=.2cm]robot.west) -- node[lab,above] {sensors} ([yshift=.2cm]driver.east);
\draw[->,thick] ([yshift=-.2cm]driver.east) -- node[lab,below] {commands} ([yshift=-.2cm]robot.west);
\draw[->,thick] (lidar) -- node[lab,above] {scan data} (lidardriver);
\node[align=center,font=\footnotesize] at (10.55,-.95) {Hardware links\\(e.g., USB serial)};
\end{tikzpicture}

\caption{A modular robot system. The five software boxes are ROS 2 nodes
in this example. Colored arrows represent topic-based data flow; black
arrows connect hardware to its drivers. Odometry feeds both mapping and
control. Start and goal inputs to planning are omitted for clarity.
The hardware can also be replaced by a simulator exposing compatible
interfaces.}
\label{fig:l14-architecture}
\end{figure}

\subsection{Nodes, topics, and messages}
A \emph{node} is a participant in the ROS 2 computation graph. It typically
has a focused responsibility, such as receiving a sensor stream or computing
control commands. The examples here run one node per executable. This is
a useful organization, not a restriction: multiple nodes can share a process,
and communicating nodes can run on different computers.

A \emph{topic} is a named communication channel. A publisher sends messages
on that topic; a subscription receives them. A node can have several
publishers and subscriptions simultaneously. Multiple subscribers can
receive the same published stream, as mapping and control do with odometry.
Topics can also have multiple publishers; they are not inherently private
point-to-point connections.

A \emph{message type} defines the structure of each message. The examples
use \texttt{std\_msgs/msg/String}, whose \texttt{data} field contains text.
The ROS message object and the Python string stored in its field are
different objects. Sensor and command interfaces use types suited to their
data rather than encoding everything as text.

For communication, the resolved topic names and message types must match,
and the quality-of-service settings must be compatible. In the default
root namespace, the relative name \texttt{my\_topic} resolves to
\texttt{/my\_topic}. A different namespace can change that resolved name.

\subsection{Why separate the components?}
If a better planner becomes available, it can replace the planner node
while leaving the mapping and control implementations intact, provided
it satisfies the same interface contract. This contract includes more
than a topic name: message structure, units, coordinate frames, meaning,
and timing requirements must agree. A path in a different frame is not
interchangeable merely because it has the same message type.

Modularity also makes it possible to inspect one stream, test one component,
or replace a real sensor with simulated data. ROS 2 supplies communication
facilities; the application still defines the algorithms and the meaning
of the exchanged information.

\section{A periodic publisher}
\label{sec:l14-publisher}
The following program, stored as \texttt{class\_20261112\_b.py}, combines
the introductory Python statements with a periodic publisher. It sends
\texttt{Ciao 0}, \texttt{Ciao 1}, and so on, with a nominal period of
$T=0.5\,\mathrm{s}$, corresponding to $f=1/T=2\,\mathrm{Hz}$.

\begin{lstlisting}[style=robotcode,language=Python]
import rclpy
from rclpy.node import Node
from std_msgs.msg import String


class MinimalPublisher(Node):
    def __init__(self):
        super().__init__('minimal_publisher')
        self.publisher_ = self.create_publisher(
            String, 'my_topic', 10)
        self.timer = self.create_timer(0.5, self.timer_callback)
        self.i = 0

    def timer_callback(self):
        msg = String()
        msg.data = 'Ciao %d' % self.i
        self.publisher_.publish(msg)
        self.i += 1
\end{lstlisting}

\begin{minipage}{\linewidth}
The remainder of the same file contains the Python demonstration and the
entry point that starts the node:
\begin{lstlisting}[style=robotcode,language=Python]
def thisIsTwooo():
    print("THIS IS TWOOOOO")


def main(args=None):
    for i in range(0, 5):
        print(i)
        if i == 2:
            print("this is two!")
    thisIsTwooo()

    rclpy.init(args=args)
    minimal_publisher = MinimalPublisher()
    try:
        rclpy.spin(minimal_publisher)
    except KeyboardInterrupt:
        pass
    finally:
        minimal_publisher.destroy_node()
        if rclpy.ok():
            rclpy.shutdown()


if __name__ == '__main__':
    main()
\end{lstlisting}
\end{minipage}

\subsection{Initialization and persistent state}
The imports make the client library \texttt{rclpy}, its \texttt{Node}
class, and the \texttt{String} message type available. The call
\texttt{rclpy.init} initializes the ROS context before the node is created.
Instantiating \texttt{MinimalPublisher} then runs its initialization method.

The name \texttt{minimal\_publisher} identifies the node in the ROS graph.
It is distinct from the Python class name and from the topic name.
The inherited method \texttt{create\_publisher} creates a publisher for
\texttt{String} messages on \texttt{my\_topic}.

The argument 10 specifies a history depth for the default keep-last QoS
configuration: a bounded number of samples can be retained. It does not
set the publication rate. Use a positive depth for this example; zero
is not a useful choice for keep-last history and can produce a warning.
Detailed QoS design is beyond this introduction.

The timer stores a reference to \texttt{self.timer\_callback}. Notice
that there are no parentheses after the callback name here: the program
passes the callable rather than executing it immediately. Finally,
\texttt{self.i = 0} initializes the counter that persists between callbacks.
The trailing underscore in \texttt{publisher\_} is simply part of the
attribute name; it has no special ROS meaning.

\subsection{Publishing a message}
Each timer callback constructs a fresh message, fills its \texttt{data}
field, publishes it, and increments the counter. In
\texttt{'Ciao \%d' \% self.i}, the integer replaces the \texttt{\%d}
placeholder. The counter is incremented after publishing, so the first
message contains zero.

Publishing does not automatically print a message in the publisher's
terminal. The Python demonstration prints once at startup; later ROS
messages are visible through a subscription or an inspection command.
Console output and topic communication are separate operations.

\subsection{Why the program needs \texttt{spin}}
Creating a timer and a publisher registers work that can happen later.
It does not by itself keep executing callbacks. The call
\texttt{rclpy.spin(node)} runs an executor that waits for ready work and
dispatches callbacks. It is an event-processing loop, not merely a sleep.
Without it, this example would create the node and then reach the end
without servicing its periodic callback.

With the default single-threaded executor, callbacks run one at a time.
A long or blocking callback can delay other work. A timer period specifies
the intended timing; it is not a guarantee of exact wall-clock execution
under every load.

The \texttt{try}/\texttt{except}/\texttt{finally} structure allows an
interactive interruption and ensures cleanup when leaving the spin call.
The node is destroyed and the context is shut down if still active.
The final \texttt{if \_\_name\_\_ == '\_\_main\_\_'} calls
\texttt{main()} when Python executes the file directly. A ROS console
entry point can instead import the module and call \texttt{main} explicitly.

\section{A subscriber and its callback}
\label{sec:l14-subscriber}
The subscriber uses the same message type and topic name. Its callback
receives a message object as an argument and prints the text through the
node's logger. Store this program as \texttt{my\_subscriber.py}.

\begin{lstlisting}[style=robotcode,language=Python]
import rclpy
from rclpy.node import Node
from std_msgs.msg import String


class MinimalSubscriber(Node):
    def __init__(self):
        super().__init__('minimal_subscriber')
        self.subscription = self.create_subscription(
            String, 'my_topic', self.listener_callback, 10)

    def listener_callback(self, msg):
        self.get_logger().info('I heard: "%s"' % msg.data)
\end{lstlisting}

\begin{minipage}{\linewidth}
The same file ends with the subscriber's entry point:
\begin{lstlisting}[style=robotcode,language=Python]
def main(args=None):
    rclpy.init(args=args)
    minimal_subscriber = MinimalSubscriber()
    try:
        rclpy.spin(minimal_subscriber)
    except KeyboardInterrupt:
        pass
    finally:
        minimal_subscriber.destroy_node()
        if rclpy.ok():
            rclpy.shutdown()


if __name__ == '__main__':
    main()
\end{lstlisting}

\end{minipage}

The subscription registers \texttt{self.listener\_callback} to process
incoming messages. The executor
invokes that callback with the received message. The callback need not
print; it could update a stored pose, process a scan, or compute a response.
The subscriber does not need its own periodic timer merely to receive data.

\begin{figure}[htbp]
\centering\begin{tikzpicture}[>=Stealth,
 box/.style={draw,rounded corners,align=center,minimum width=3.4cm,minimum height=1.5cm}]
\node[box] (pub) at (0,0) {\texttt{minimal\_publisher}\\\texttt{timer\_callback()}};
\node[box] (sub) at (8.5,0) {\texttt{minimal\_subscriber}\\\texttt{listener\_callback(msg)}};
\draw[->,thick,sensorframe] (pub) -- node[above,font=\footnotesize] {\texttt{/my\_topic}}
 node[below,font=\footnotesize] {\texttt{String: data="Ciao 0"}} (sub);
\draw[->] (0,1.8) node[above] {timer: $T=0.5\,\mathrm{s}$} -- (pub.north);
\draw[->] (8.5,1.8) node[above] {message received} -- (sub.north);
\node[font=\footnotesize] at (0,-1.15) {executor processes ready callbacks};
\node[font=\footnotesize] at (8.5,-1.15) {executor processes ready callbacks};
\end{tikzpicture}

\caption{Two different events drive the two callbacks. A timer makes the
publisher callback ready; a received message makes the subscriber callback
ready. Each node's executor dispatches its callbacks while spinning.
The topic carries messages, not Python function calls between the programs.}
\label{fig:l14-pubsub}
\end{figure}

A subscriber started after the publisher normally begins with newly
received messages rather than the complete earlier sequence under these
default settings. Likewise, a history depth of 10 does not mean that a
late subscriber automatically receives the last ten messages.

\section{Packages, entry points, and execution}
\label{sec:l14-running}
A workspace groups packages and their build products. A Python ROS package
contains both Python modules and metadata used to install and locate
executables.

\begin{minipage}{\linewidth}
The relevant structure in this example is
\begin{lstlisting}[style=robotcode]
ros2_ws/
  src/
    my_package/
      package.xml
      setup.py
      setup.cfg
      resource/my_package
      my_package/
        __init__.py
        class_20261112_b.py
        my_subscriber.py
\end{lstlisting}
\end{minipage}

The outer \texttt{my\_package} is the ROS package directory; the inner
directory is its importable Python package. The workspace's \texttt{build},
\texttt{install}, and \texttt{log} directories are generated products.
For navigation, \texttt{pwd} prints the working directory, \texttt{ls}
lists its contents, \texttt{cd} changes directory, and \texttt{cd ..}
moves to the parent directory.

\subsection{Registering the executable}
Saving a Python file does not register it with \texttt{ros2 run}.
In \texttt{setup.py}, console-script entries map executable names to
Python functions. The two
relevant entries, within the existing \texttt{setup(...)} call, are
\begin{lstlisting}[style=robotcode,language=Python]
entry_points={
    'console_scripts': [
        'class_20261112_b = my_package.class_20261112_b:main',
        'listener = my_package.my_subscriber:main',
    ],
},
\end{lstlisting}
Each entry has the form \texttt{executable = package.module:function}.
For example, \texttt{listener} is the executable name; it imports
\texttt{my\_package.my\_subscriber} and calls \texttt{main}. Add these
entries alongside other executables rather than removing unrelated entries.
The package declares dependencies including \texttt{rclpy} and
\texttt{std\_msgs} in \texttt{package.xml}; \texttt{setup.cfg} places
the installed scripts where ROS tooling expects to find them.

\subsection{Build, source, and run}
The commands below assume a Linux shell with the ROS 2 installation already
available, such as a prepared container. From the workspace root:
\begin{lstlisting}[style=robotcode]
cd ~/ros2_ws
colcon build --packages-select my_package
source install/setup.bash
ros2 run my_package class_20261112_b
\end{lstlisting}
For a Python package this build step primarily installs code and entry
points; it is not C-style compilation. With an ordinary copied installation,
save and rebuild after edits. A workspace configured with symlink installation
can expose many Python source edits directly, but metadata changes still
require rebuilding. Sourcing the workspace makes its installed packages
available in that shell. A new terminal needs the environment too.

In a second terminal, with the base ROS environment available:
\begin{lstlisting}[style=robotcode]
source ~/ros2_ws/install/setup.bash
ros2 run my_package listener
\end{lstlisting}
The subscriber should report messages such as \texttt{I heard: "Ciao 12"}.
The initial number depends on when it connects. The executable can be run
from any directory once its environment is sourced; the build commands
above deliberately use the workspace root.

\subsection{Observing the communication}
The command-line tools can inspect the running system without adding
application code:
\begin{lstlisting}[style=robotcode]
ros2 node list
ros2 topic list
ros2 topic info /my_topic
ros2 topic echo /my_topic
\end{lstlisting}
The echo command acts as another subscriber and displays messages as they
arrive. It can run alongside the custom subscriber. Tab completion can
help discover available package and executable names.

If no messages arrive, first check that both processes are running and
spinning, then inspect the resolved topic names, types, and QoS compatibility.
In particular, \texttt{topic} and \texttt{my\_topic} are different names.
If a new executable is not found, check its entry point, rebuild, and
source the correct workspace. Distinguish the executable name, ROS node
name, and topic name: they serve different purposes and need not match.

\clearpage
\section{Additional exercises}
\begin{enumerate}
\item \textbf{Indentation changes behavior.} Write the output of the
introductory Python example for each of the three placements of
\texttt{thisIsTwooo()}: outside the loop, inside the loop but outside
the conditional, and inside the conditional. Explain the order of the output.
\item \textbf{Persistent state.} Predict what happens if \texttt{self.i = 0}
is moved from the publisher's initialization method to the beginning of
\texttt{timer\_callback}. Why is instance state useful here?
\item \textbf{Two timing parameters.} Change the timer period to
$0.2\,\mathrm{s}$. What is the nominal publication frequency? Would
changing the history depth from 10 to 20 have the same effect?
\item \textbf{A modified subscriber.} Change the callback so that it
prints only messages whose counter is even. State how you extract the
counter from the received text and test the first six messages.
\item \textbf{Independent receivers.} Run the custom subscriber and
\texttt{ros2 topic echo /my\_topic} simultaneously. Predict what happens
when either receiver is stopped. Does the publisher require both to exist?
\item \textbf{Find the mismatch.} A publisher uses \texttt{topic} and
a subscriber uses \texttt{my\_topic}, both with \texttt{String} messages
in the root namespace. Explain why they do not exchange messages. Identify
which inspection commands would help find the problem.
\item \textbf{A replaceable planner.} Using
Figure~\ref{fig:l14-architecture}, list the inputs and outputs of a planner
that could replace the existing one. Besides the message types, give three
interface properties that must remain compatible. Explain why changing
the planning algorithm need not change the lidar driver.
\end{enumerate}

\smallskip
\textit{Further reading.} The official ROS 2 tutorial
\href{https://docs.ros.org/en/humble/Tutorials/Beginner-Client-Libraries/Writing-A-Simple-Py-Publisher-And-Subscriber.html}{\emph{Writing a simple publisher and subscriber (Python)}}
provides a related example of this programming pattern.
